%% implementation.tex
%%

%% ==============================
\section{Implementation}
\label{sec:Implementation}
%% ==============================

This section describes the final implementation of the dynamic Multi-Armed Bandit (MAB) hyper-heuristic system. It covers the core software stack, inter-process orchestration pipeline, low-level heuristic pool, and structured data outputs.

\subsection{Technology Stack and Software Architecture}
The framework uses a two-tier software architecture built on two primary programming environments:
\begin{itemize}
    \item \textbf{Python Execution Engine (v3.11.7):} Serves as the primary control plane; housing the multi-armed bandit decision algorithms, workload trace ingestion module, statistical workload fingerprinter and reward transformer. It uses \texttt{NumPy} and \texttt{SciPy} for statistical operations and \texttt{Pandas} for trace parsing.
    \item \textbf{Java Simulation Engine (OpenJDK 21 / CloudSim Plus v7.3.0):} Acts as the cloud simulation environment. CloudSim Plus core classes were extended to support low-level heuristic instantiation and I/O streams
\end{itemize}

\subsection{Standard I/O Inter-Process Communication Pipeline}
In order to reduce the overhead of the spawning a new process for every CloudSim instance, we make use of inter-process communication to create a singular data stream between the Python control plane and the Java simulation engine.

The Python program launches the Java instance as a managed subprocess that receives a JSON payload containing the batch index and task specifications.

The Java application holds a stream reader that parses incoming JSON records asynchronously, which then executes the batch across the virtual machine pool.

\subsection{Low-Level Heuristic Pool Implementation}
The low-level heuristic pool is implemented inside the Java simulation core under a unified strategy interface, with nine distinct scheduling routines being implemented:
\begin{enumerate}
    \item \textbf{List-Based Schedulers (HEFT, PEFT, Min-Min, Max-Min, Sufferage, HEFT-TD):} Custom task prioritization routines construct directed acyclic graph (DAG) task ranks and evaluate expected finish times across heterogeneous VM capacity matrices on the fly.
    \item \textbf{Deterministic Baseline (Round-Robin):} Sequentially maps incoming cloudlets to active VMs to establish a dynamic baseline for performance normalization.
    \item \textbf{Local Search and Meta-Heuristics (Simulated Annealing, HGHHC):} Algorithmic state models execute local schedule permutations. HGHHC integrates Henry Gas Solubility Optimization and Harris Hawks Optimization search mechanics, running iterative candidate schedule evaluations within bounded CloudSim time slots.
\end{enumerate}

\subsection{Energy, Carbon Footprint, and Cost Modeling}
To evaluate multi-objective trade-offs accurately, custom tracking modules were integrated directly into the CloudSim Plus Datacenter engine:
\begin{itemize}
    \item \textbf{Monetary Cost Model:} This models dynamic cloud resource pricing, with the help of AWS billing API. Each VM instance incurs a fixed per-second compute rate based on its allocated processing capacity (MIPS) and memory profile, following tier-based provider pricing schemes.
    \item \textbf{Dynamic Power Model:} Power consumption $P(u)$ is calculated as a function of instantaneous CPU utilization $u \in [0, 1]$:
    \begin{equation}
    P(u) = P_{\text{idle}} + (P_{\text{max}} - P_{\text{idle}}) \times u
    \end{equation}
    where $P_{\text{idle}}$ represents baseline VM standby power, and $P_{\text{max}}$ denotes peak power under full load.
\end{itemize}

\subsection{Non-Stationary Infrastructure Fault Injection}
An automated fault injection module was implemented in the Python environment to simulate physical infrastructure failure. The purpose of this is to monitor and evaluate how each bandit algorithm reacts to an instantaneous change in physical environments. The module shuts down seven of the eight provisioned virtual machines, thereby forcing all incoming tasks to a single machine.

\subsection{Generated Outputs and Data Transformation}
The execution pipeline produces artefacts for post-hoc analysis:
\begin{enumerate}
    \item \textbf{Raw Trace Log Stream:} JSON files generated per run that hold batch-level decision metrics such as selected heuristic arm, raw makespan, monetary cost, carbon output, VM utilization, and squashed reward $R_k$.
    \item \textbf{Aggregated Seed Benchmark Dataset:} Structured CSV files combining metrics across all 10 random seeds per MAB variant and trace dataset.
\end{enumerate}
